\section{Model Design and Implementation}
\label{ch:model}

We specify the \ac{ABM} following the \ac{ODD} protocol \cite{grimm2006odd,
grimm2020odd}. Per the protocol's \emph{purpose and patterns} framework, empirical target criteria are fixed prior to presenting simulation outputs.
Tables~\ref{tab:design-concepts} and~\ref{tab:submodels} provide the formal model specification; the accompanying text details the core decision rules and transmission submodels.
The model is given as pseudocode in Appendices~\ref{app:pseudo-init} and~\ref{app:pseudo-main}, which are the normative statement of every rule below; rejected structural alternatives and supplementary material appear in Appendices~\ref{app:rationale} and~\ref{app:supplementary}.

\subsection{Purpose and Patterns}
\label{sec:purpose}

The simulation employs a counterfactual injection design.
Each household initializes with a zero \ac{BNPL} balance, the macroeconomic environment is fixed at 2017 baseline conditions, and the reference model is calibrated with \ac{BNPL} disabled.
\ac{BNPL} access is subsequently introduced while holding all other parameters constant, isolating the systemic effect of the added credit channel.
Consequently, simulation outputs do not constitute macro forecasts for a specific calendar year; rather, they quantify the marginal impact of \ac{BNPL} availability.

Table~\ref{tab:patterns} lists four empirical patterns established prior to simulation execution. Pattern~2, the 2017-Q1 National Credit Regulator (\ac{CCMR}) arrears profile (Table~\ref{tab:ccmr}), serves as the primary calibration target.
Two model parameters are calibrated to match its 90+ day and 1--30 day default bands, leaving intermediate bands unconstrained.
Patterns 1, 3, and 4 represent out-of-sample or off-vintage benchmarks used as secondary plausibility checks, for the reasons given below.
Model performance against the unconstrained \ac{CCMR} intermediate bands thus provides the primary out-of-sample validation for the first research question.

\label{sec:lim-validation}
The four patterns are not of equal standing.
Only Pattern~2 is South African and vintage-matched, and it is measured on accounts while the model's unit is the household, so the comparison asks only whether the arrears profile has the right shape and scale, and that imprecision propagates into the fitted shock probability; arrears proportions are accordingly computed over the credit-active subpopulation (the 44\% of households holding traditional debt), with the population-wide default rate tracked separately.
Pattern~1 is partly designed in: the want-driven borrowing path of Submodel~3 was adopted because a purely shortfall-driven agent could not reproduce complementarity under any parameterisation, so its direction is built in while its magnitude is not.
Patterns~3 and~4 rest on a United Kingdom model and a 2025 survey respectively.
All three are reported as plausibility checks, not tests.

\begin{table}[H]
\centering
\caption{Empirical patterns registered before any result was produced}
\label{tab:patterns}
\tabnote{The four empirical patterns fixed before simulation, following the pattern-oriented modelling convention of the \ac{ODD} protocol. Pattern~2 is the calibration target: two parameters, the income-shock probability and the payment-friction rate, are fitted to its 90+ and 1--30 day bands, so agreement on those bands is not evidence. Patterns~1, 3 and 4 are foreign or off-vintage and are read as plausibility checks, as argued in the text. The third column states what the model must produce for the pattern to count as reproduced and what would falsify it.}
\begin{tabular}{L{2.4cm}L{4.2cm}L{4.8cm}L{1.5cm}}
\toprule
\textbf{Pattern} & \textbf{Empirical Observation} & \textbf{Model Analogue \& Falsification Condition} & \textbf{Source} \\
\midrule
1. Complementarity, not substitution & \ac{BNPL} adoption correlates with subsequent increases in credit-card interest and late fees. & Enabling \ac{BNPL} should elevate arrears and servicing stress on \emph{traditional} debt. A reduction falsifies the lender-choice hierarchy. Measured via servicing stress, not debt stock. & \cite{dehaan2024bnpl} \\
\addlinespace
2. Baseline arrears & 2017-Q1 arrears profile for unsecured credit and credit facilities combined. & With \ac{BNPL} disabled, the credit-active population's arrears distribution must approximate Table~\ref{tab:ccmr}. \textbf{(Calibration target)} & \cite{ncr_ccmr_2017} \\
\addlinespace
3. Non-monotonic debt burden & Debt-to-income peaks among middle-income households in independent modelling. & The aggregate debt-to-income ratio (total debt over total income per quintile, the pre-registered statistic) must peak mid-distribution rather than rising monotonically with income. & \cite{hamill2023creditcard} \\
\addlinespace
4. Savings exhaustion & 36\% of South African consumers anticipated missing a bill payment. & The share of agents reaching zero liquid savings should be of the same order. A stated expectation compared against a realised stock, so order-of-magnitude only. & \cite{transunion_cps_sa_2025} \\
\bottomrule
\end{tabular}
\end{table}

\begin{table}[H]
\centering
\caption{2017-Q1 \ac{CCMR} arrears profile}
\label{tab:ccmr}
\tabnote{The calibration target: the distribution of accounts by days in arrears for unsecured credit and credit facilities combined, first quarter of 2017, from Tables 21 and 23 of the \ac{CCMR} \cite{ncr_ccmr_2017}. The unit is the account, not the household, so the model's household-level arrears distribution over the credit-active subpopulation is compared with it for shape and scale rather than point for point. The two aggregate rows are cumulative over the bands above them.}
\begin{tabular}{lr}
\toprule
\textbf{Age band} & \textbf{Share of accounts (\%)} \\
\midrule
Current                  & 71.63 \\
1--30 days               &  8.24 \\
31--60 days              &  3.59 \\
61--90 days              &  2.32 \\
91--120 days             &  1.76 \\
120+ days                & 12.45 \\
\midrule
\emph{Aggregate: 60+ days}  & \emph{16.54} \\
\emph{Aggregate: 90+ days}  & \emph{14.21} \\
\bottomrule
\end{tabular}
\end{table}

\subsection{Entities, Scales and Scheduling}
\label{sec:entities}

The model consists of five entities:
\begin{itemize}
    \item \textbf{Household agents} are the adaptive entities. Each maps to an empirical
    \ac{NIDS} Wave 5 household (Table~\ref{tab:agent-mapping}) and tracks cash reserves,
    traditional debt, \ac{BNPL} obligations, arrears, distress, default status, and a fixed
    repayment archetype (minimum-payer or scheduled-payer).
    \item \textbf{The traditional lender} is a single aggregate entity holding all traditional
    debt. It is non-adaptive and acts only as a deterministic affordability gate.
    \item \textbf{The credit bureau} is a shared registry of traditional debt and defaults that
    defines what the traditional lender can observe.
    \item \textbf{\ac{BNPL} platforms} are the injected entities, four in the baseline to mirror
    the South African market. Each holds isolated receivables, applies an automated screen with a
    per-order cap and per-household rolling limits, and has no visibility into competitors.
    \item \textbf{Reference groups} define the social topology: each household belongs to one
    group, its income quintile crossed with its province.
\end{itemize}

Figure~\ref{fig:visibility} illustrates the model's information structure.
The dashed connections highlight the reporting gaps of the pre-2026 position: platforms lack visibility into competitor extensions, and \ac{BNPL} liabilities are not reported to the central bureau.

\begin{figure}[H]
\centering
\begin{tikzpicture}[
    font=\small,
    box/.style={draw, align=center, inner sep=5pt, rounded corners=1pt},
    sees/.style={->, >=Stealth, semithick},
    gap/.style={->, >=Stealth, semithick, dashed},
    lbl/.style={font=\footnotesize, align=center},
]
\node[box, minimum width=4.0cm] (hh) at (0,0) {\textbf{Household}\\[1pt]
    \footnotesize income, savings, traditional debt,\\
    \footnotesize \ac{BNPL} obligations, informal debt};
\node[box] (lender) at (0,3.0) {\textbf{Traditional lender}\\ \footnotesize Regulation 23A gate};
\node[box] (bureau) at (7.0,3.0) {\textbf{Credit bureau}\\ \footnotesize traditional debt, defaults};
\node[box] (p1) at (6.0,0.35)  {\footnotesize Platform 1};
\node[box] (p2) at (8.2,0.35)  {\footnotesize Platform 2};
\node[box] (p3) at (6.0,-0.75) {\footnotesize Platform 3};
\node[box] (p4) at (8.2,-0.75) {\footnotesize Platform 4};
\node[draw, dashed, inner sep=6pt, fit=(p1)(p2)(p3)(p4),
      label={[lbl]below:{mutually blind (Submodel 13)}}] (plats) {};
\node[box] (group) at (0,-2.6) {\textbf{Reference group}\\ \footnotesize income quintile $\times$ province};

\draw[sees] (hh) -- (lender) node[lbl, midway, left] {applications:\\ declared income};
\draw[sees, <->] (lender) -- (bureau) node[lbl, midway, above] {reports and reads};
\draw[sees] (hh.east) -- (plats.west) node[lbl, midway, above] {pay-in-four\\ facilities};
\draw[gap] (plats.north) -- (bureau.south) node[lbl, midway, right] {no reporting duty\\ (Scenario 1 creates it)};
\draw[sees] (group) -- (hh) node[lbl, midway, left] {lagged adoption\\ share $s_g$};
\end{tikzpicture}
\caption{Information structure of the model}
\label{fig:visibility}
\tabnote{Boxes are the five entity types; solid arrows are information flows that exist in every run; dashed connections are the two information gaps of the pre-2026 South African position. Platforms cannot see one another's exposures (Submodel~13), and \ac{BNPL} originations do not reach the bureau unless the reporting switch of Submodel~15 (Scenario~1, the 2026 requirement) is on. The traditional lender therefore evaluates the Regulation 23A test against a bureau file that omits \ac{BNPL}, and households observe only their own accounts and the lagged adoption share $s_g$ of their reference group (Submodel~11).}
\end{figure}

The simulation updates across discrete 14-day ticks, mapping standard biweekly pay-in-four installment schedules directly to tick steps.
Monthly survey cash flows are scaled to ticks using a factor of $12/26$.
Simulation runs span 52 ticks (24 calendar months), discarding the first 12 ticks as initialization burn-in.
Within each tick, schedule execution follows a synchronous sequence.
Household activation order is randomized uniformly at each tick, and credit decisions execute dynamically as applications occur.
Social interaction signals update asynchronously at the end of each tick step.
A household's tick proceeds in seven steps, given in full as Algorithm~\ref{alg:tick}: income arrives, subject to the employment shock; committed expenditure is paid; interest accrues and scheduled service falls due, the contractual instalment for scheduled-payers or interest plus a minimum principal for minimum-payers, together with \ac{BNPL} instalments and arrears; a realised shortfall triggers a credit application, \ac{BNPL} first where the household is eligible and the traditional lender for the remainder, and distress is assessed on what remains unmet; \ac{BNPL} debits are collected and traditional service is paid, subject to payment friction; discretionary consumption is set from the remainder; and a want-driven \ac{BNPL} purchase is drawn before balances, arrears age, distress and default update.

Three sequence choices govern transmission dynamics within the model.
First, prioritizing committed consumption over debt service enables cash-flow default driven by living expense shortfalls \cite{madeira2018chile}.
Second, allowing credit applications prior to debt settlement models debt-servicing liquidity feedback \cite{dorazio2017micro}.
Third, evaluating distress prior to payment friction ensures that inadvertent payment delays do not misclassify solvent households as insolvent.
Lagging the reference group signal by one tick ensures that peer influences remain independent of intra-tick execution order.

\subsection{Design Concepts}
\label{sec:design-concepts}

Table~\ref{tab:design-concepts} outlines the eleven structural \ac{ODD} design concepts, and closes with the six mechanisms the model excludes by design.

\begin{longtable}{L{2.9cm}L{7.3cm}L{3.2cm}}
\caption{\ac{ODD} design concepts and design exclusions}
\label{tab:design-concepts} \\
\multicolumn{3}{p{13.4cm}}{\footnotesize The eleven design concepts of the \ac{ODD} protocol \cite{grimm2020odd}, each with the model's treatment and the literature that grounds it; ``None'' in the source column marks a modelling choice with no external anchor. The final row lists the six mechanisms excluded by design and the reason for each. The second exclusion, the absence of any channel by which one household's default affects another, is the structural explanation of the central negative result of Section~\ref{sec:results-rq2}.} \\
\toprule
\textbf{Concept} & \textbf{Treatment} & \textbf{Source} \\
\midrule
\endfirsthead
\toprule
\textbf{Concept} & \textbf{Treatment} & \textbf{Source} \\
\midrule
\endhead
\bottomrule
\endfoot
Basic principles & Balance sheets form the accounting core. Heuristic decision rules replace intertemporal optimisation. Statutory regulations act as explicit computational gates. & \cite{cardaci2018inequality, dorazio2017micro, meier2010present, Keys2019, yun2020housing} \\
\addlinespace
Emergence & Default rates, stacking depths and adoption paths are not imposed. They emerge from two separable feedback loops: a \emph{financial} loop, in which borrowing to service existing debt inflates future obligations, and a \emph{social} loop, in which lagged group adoption raises individual propensities. $\beta=0$ disables the social loop only, so the two are separately attributable. & None \\
\addlinespace
Adaptation & Households compress consumption and select borrowing and repayment actions from immediate liquidity alone. The traditional lender is non-adaptive. & \cite{dorazio2017micro} \\
\addlinespace
Objectives & No forward-looking utility. Households satisfy obligations in a fixed priority order, compressing discretionary spending under stress. & \cite{meier2010present} \\
\addlinespace
Learning & Excluded. Behavioural rules are static over the horizon. & None \\
\addlinespace
Prediction & Excluded. Agents treat current income as persistent. & None \\
\addlinespace
Sensing & Agents know their own balance sheets and observe the lagged adoption share among their reference group's \ac{BNPL}-eligible members. The traditional lender is blind to \ac{BNPL} and informal debt. & \cite{nortonrose_bnpl_sa} \\
\addlinespace
Interaction & Households interact with lenders through credit applications, and with each other indirectly through reference-group adoption shares. & \cite{ackert2025bnpl} \\
\pagebreak
Stochasticity & Income shocks, archetype and eligibility assignment, purchase-size draws, platform routing order, payment friction, and threshold endowments (drawn from a separate initialisation stream so control arms reproduce exactly). Runs are seeded and replicated, with dispersion reported. & None \\
\addlinespace
Collectives & Reference groups (income quintile $\times$ province) aggregate individual adoption into the peer signal. & \cite{cardaci2018inequality} \\
\addlinespace
Observation & Tracked per tick: population default rate, \ac{CCMR} arrears distribution, aggregate traditional and \ac{BNPL} debt volumes, mean and quintile \ac{DSTI}, consumption compression, and application and rejection counts. & None \\
\addlinespace
Exclusions & Six mechanisms are excluded by design. \emph{No macroeconomic dynamics}: any time-variation would confound the injection effect with a macro effect. \emph{No contagion}: one aggregate lender with no balance sheet that can fail, and no feedback from household default to income, so defaults are correlated individual events and cannot cascade. \emph{No competition among traditional lenders}, which suppresses the lender-conduct mechanism of Hamill et al. \emph{No learning and no scarring}. \emph{No peer effects on consumption}: peer influence acts only on \ac{BNPL} adoption. \emph{Static household composition}. & \cite{hamill2023creditcard, cardaci2018inequality} \\
\end{longtable}

\subsection{Submodels}
\label{sec:submodels}

Initial household states are specified using the empirical microdata from Section~\ref{ch:population}, setting initial \ac{BNPL} balances to zero.
Table~\ref{tab:submodels} details the seventeen operational submodels; key structural mechanisms are elaborated below.

\begingroup
\setlength{\tabcolsep}{4pt}
\small
\begin{longtable}{L{2.35cm}L{4.75cm}L{2.3cm}L{4.0cm}}
\caption{Submodel specification}
\label{tab:submodels} \\
\multicolumn{4}{p{13.4cm}}{\footnotesize The seventeen submodels, each with its rule, the source grounding it, and its parameters. Every behavioural rule is grounded in literature or flagged as an assumption; every free parameter is sourced, derived or swept, in the sense defined in the caption of Table~\ref{tab:param-register}, which lists every parameter with its baseline value, provenance and sweep range. Bold marks the rules with no literature anchor, for which sensitivity analysis is mandatory. Submodels 10, 11, 13 and 15 are elaborated in the text; the rest are given in full as pseudocode in Appendix~\ref{app:pseudo-main}.} \\
\toprule
\textbf{Submodel} & \textbf{Rule} & \textbf{Source} & \textbf{Parameters and validation} \\
\midrule
\endfirsthead
\toprule
\textbf{Submodel} & \textbf{Rule} & \textbf{Source} & \textbf{Parameters and validation} \\
\midrule
\endhead
\bottomrule
\endfoot

1. Income and shock & Baseline income is static. Each employed member of a wage-dominant household separates with a per-tick probability, and the household loses that member's share of its wage component until re-employment. Grant- and remittance-dominant households are exempt: grants are statutory transfers. & \cite{madeira2018chile, statssa_lmd_2022} & Onset probability fitted to the \ac{CCMR} 90+ band and reported against the \ac{QLFS} 2017 separation band of 0.54--1.17\% per tick. No magnitude parameter: the loss is the observed wage component. Exit hazard 1.87\% per tick from \ac{QLFS} 2017 Q3--Q4 flows. Single-tick non-persistent shock kept as a robustness arm. \\
\addlinespace
2. Consumption & Two tiers. Committed expenditure (food, rent) is an incompressible floor; discretionary spending compresses to zero before any instalment is missed. & Data layer (Section~\ref{ch:population}) & Baseline floor of zero; habit-persistence floors of 0.25 and 0.50 swept. \\
\addlinespace
3. Borrowing trigger & Traditional credit is sought only for realised shortfalls. \ac{BNPL} adds a want-driven path governed by peer adoption. & \cite{dorazio2017micro, ackert2025bnpl, dehaan2024bnpl} & $q_{\text{base}}$, $\beta$; see Submodel 11. \\
\addlinespace
4. Borrowing amount & Shortfall path: \textbf{structural assumption} with no anchor located; households request the exact shortfall. Want-driven path: purchase size is drawn lognormally with mean $\kappa$ times the household's own monthly discretionary budget, so purchase size inherits the population's heterogeneity. & \cite{statssa_ies_2023} & \textbf{Mandatory sweeps:} shortfall / $+25\%$ / $+$ one tick of committed expenditure; $\kappa = 0.139$ (derived from IES 2022/23 COICOP shares) swept 0.07--0.28; dispersion 0.6 (assumption); sizing base swept discretionary / income. Realised mean order checked against the R992 (2017 Rands) issuer-disclosed average order value \cite{weaver2025iar}. \\
\addlinespace
5. Lender choice & Eligible households take \ac{BNPL} first; any remainder routes to traditional credit. The baseline uses traditional credit only. & \cite{ackert2025bnpl} & No free parameter. Falsification tied to Pattern 1. \\
\addlinespace
6. Repayment and arrears & Scheduled-payers pay the full amortised instalment. Minimum-payers pay the larger of accrued interest and 5\% of balance per month, tick-scaled; the fraction is an \textbf{assumption}, since neither the \ac{NCA} nor any located source prescribes a formula. Unpaid amounts accrue as arrears in \ac{CCMR} age bands. Payment friction skips an affordable traditional instalment with fixed probability. & \cite{Keys2019, Kuchler2021} & Minimum-payer share 0.29 (swept 0.20--0.40); minimum fraction swept 0.025--0.10. Friction is the second fitted parameter, calibrated to the \ac{CCMR} 1--30 day band; it applies to traditional debt only, since \ac{BNPL} auto-debits a card. \\
\addlinespace
7. Distress and default & A household is distressed when income, savings and any credit raised cannot cover committed expenditure plus scheduled service, assessed before payment friction. Seven consecutive distressed ticks trigger default; a defaulted household is refused all further credit. & \cite{madeira2018chile} & $k=7$ ticks (98 days), the first tick boundary at or past the 90-day impairment convention; $k=4$ (60+ band) swept. \\
\addlinespace
8. Heterogeneity & Behavioural parameters vary by income quintile, demographics, and the financial-inclusion markers of the data layer. & \cite{Keys2019} & No continuous present-bias parameter: no local calibration data. \\
\addlinespace
9. Credit-granting gate & The Regulation 23A residual-income test, applied at household level; all-or-nothing. New credit is priced as unsecured (repo $+21\%$, 25-month term), matching Section~\ref{sec:servicing}. & \cite{sa_ncr_affordability_2015} & No free parameter; the statutory expense table is hard-coded and asserted against the regulation's published worked examples. \\
\addlinespace
10. Information asymmetry & The bureau records traditional debt and defaults only. Lenders cannot observe \ac{BNPL} or informal debt. & \cite{nortonrose_bnpl_sa, transunion_cps_sa_2025} & \texttt{bnpl\_bureau\_visible} boolean (default off, the pre-2026 position); Scenario 1 in Submodel 15. \\
\addlinespace
11. Peer influence & $q_i(t) = \mathrm{clip}(q_{\text{base}} + \beta s_{g(i)}(t-1), 0, 1)$ within reference groups, where $s_g$ is measured on the group's \ac{BNPL}-eligible members. A heterogeneous-threshold variant is pre-registered as the structural alternative and run on identical grids. & \cite{ackert2025bnpl, cardaci2018inequality, granovetter1978threshold} & \textbf{$\beta=0$ control arm required in all comparisons}; the threshold arm reports its own $\gamma=0$ row. \\
\addlinespace
12. \ac{BNPL} eligibility & Banked households only; a light automated screen, deliberately not Regulation 23A. Per-order cap ZAR 9{,}927 (2017 Rands); rolling limit $\lambda = 0.10$ of monthly income per platform, set per household per the industry's disclosed ``low and grow'' practice. & \cite{payflex_terms, payjustnow_terms, homechoice2024ir} & Access capped at the empirical 82.8\% banked share. The 0.10 baseline is set so that four per-platform limits sum to the stacked total Woolard observed, roughly 37\% of monthly income across providers; $\lambda$ is swept 0.1--1.0 and reported against a 0.10--0.26 band converted by the author from UK and Australian disclosures \cite{woolard2021, asic2020rep672}. Order-cap and limit binding rates reported by quintile. \\
\addlinespace
13. Stacking & Concurrent facilities across platforms, with no aggregate limit. Platforms stay blind to competitor exposures. & Follows Submodel 10 & Platforms swept 1--6 (baseline 4); emergent depth benchmarked against the \ac{CFPB} shares of Section~\ref{sec:background}. \\
\addlinespace
14. \ac{BNPL} repayment & Exogenous ``Pay-in-4'': 25\% at checkout, then 25\% at each of the next three ticks, interest-free. The checkout debit must clear, so an order is truncated to what the household can fund and liquidity caps purchase size endogenously. Late fees of ZAR 125.74 per tick, capped at ZAR 188.61 per facility (2017 Rands); a household that exhausts the cap is cut off by that platform only. & \cite{payflex_terms} & No free parameters; instalments fall on tick boundaries. Pay-in-3 monthly (PayJustNow) is the structural alternative arm. \\
\addlinespace
15. Scenario switches & Bureau visibility ($\mathbb{1}^{\text{bur}}$: \ac{BNPL} obligations enter the Regulation 23A test), mandatory affordability screening ($\mathbb{1}^{\text{aff}}$: the Regulation 23A test applied to \ac{BNPL} originations), and peer influence ($\beta > 0$), crossed. & \cite{businessday_bnpl_2026, fca_ps26_1, woolard2021} & Binary switches; every scenario is reported at $\beta = 0$ and $\beta > 0$. A cooling-off window and a concurrent-facility cap are implemented but reported only in Appendix~\ref{app:supplementary}. \\
\addlinespace
16. Macro environment & Static. The repurchase rate is fixed at 7.00\%, freezing statutory maximum interest rates. & Experimental design (Section~\ref{sec:purpose}) & Exogenous constraint. \\
\addlinespace
17. Activation order & Random, re-drawn each tick; applications decided inline as each household acts. & \cite{comer2013activation, alizadeh2015activation} & Uniform and synchronous robustness arms (Appendix~\ref{app:supplementary}). \\
\end{longtable}
\endgroup

\subsubsection*{Submodel 4: Borrowing Amount}
\label{sec:amount}

Submodel~4 governs how much a household asks for when it comes up short, and no located study settles it.
It is the model's first uncited rule, so all three specifications in Table~\ref{tab:submodels} are run, and Section~\ref{sec:results-robustness} reports the rule as the largest bar on the tornado.
Two of the three agree to within the replicate noise floor; the range comes from the third and most generous, so the sensitivity reflects one unconstrained variant rather than general instability.

\subsubsection*{Submodels 6 and 8: Imported Behavioural Parameters}
\label{sec:lim-transfer}

The behavioural micro-foundations of Section~\ref{sec:lit-behavioural} are United States evidence applied to South African borrowers: the minimum-payer share of 0.29, the present-bias results behind payment friction, and the experiment motivating the want-driven path.
No developing-market evidence of comparable quality was located for any of them.
The share that matters most is swept from 0.20 to 0.40, and the mechanisms are claims about decision-making under credit contracts rather than about United States institutions, but neither argument establishes that the parameters are right for South Africa; for the same reason Submodel~8 carries no continuous present-bias parameter, since no South African distribution exists to calibrate one.

\subsubsection*{Submodel 10: Information Asymmetry}
\label{sec:asymmetry}

The traditional lender assesses credit applications using self-reported gross income, existing internal exposures, and active traditional balances logged at the central bureau.
Unregulated liabilities, such as \ac{BNPL} commitments \cite{nortonrose_bnpl_sa, transunion_cps_sa_2025} and informal loans (e.g., stokvels or mashonisa credit) \cite{finscope2019}, remain unobserved. Consequently, statutory Regulation 23A affordability checks are applied to incomplete liability profiles, capturing a key regulatory oversight gap rather than lender non-compliance.

Bureau reporting is toggled via a boolean parameter (\texttt{bnpl\_bureau\_visible}), set to false in baseline runs to reflect the South African position before the 2026 reporting requirement.
Enabling bureau visibility incorporates active \ac{BNPL} obligations into formal affordability checks, isolating the systemic impact of unobserved credit; it is the model's representation of the requirement now being implemented (Section~\ref{sec:background}).

That representation is partial.
The switch makes \ac{BNPL} obligations visible to the lender's residual-income calculation and to nothing else: the model carries no credit score, so it cannot represent the second channel through which reported \ac{BNPL} data will act, a change in consumers' scores under the bureaus' models.
That channel is not minor: one bureau's preliminary analysis suggests the reported data could lower the scores of roughly a quarter of credit-active consumers under current models \cite{finasa_bnpl_2026}, which would ration credit through pricing and refusal rather than through affordability arithmetic.
The bureau-visibility contrast of Section~\ref{sec:results-rq3} is therefore a statement about the affordability channel only; its near-null result is not evidence that the requirement will have no effect, and a scoring rule is the extension the 2026 change most directly motivates (Section~\ref{sec:future-work}).

\subsubsection*{Submodel 11: Peer Influence}
\label{sec:peer-influence}

Let $s_g(t-1)$ represent the proportion of \ac{BNPL}-eligible households in reference group $g$ holding active \ac{BNPL} balances at tick $t-1$. Restricting the calculation denominator to eligible group members prevents adoption signals from scaling directly with overall financial inclusion rates. Household $i$'s discretionary origination propensity is defined by:

\begin{equation}
    q_i(t) = \mathrm{clip}\!\left(q_{\text{base}} + \beta \, s_{g(i)}(t-1),\; 0,\; 1\right).
    \label{eq:peer}
\end{equation}

A non-zero baseline propensity ($q_{\text{base}} > 0$) ensures non-zero adoption at initialization when aggregate balances are zero.
The parameters $q_{\text{base}}$ and $\beta$ govern baseline innovation and peer imitation dynamics, analogous to a Bass diffusion process.
Reflecting findings that \ac{BNPL} adoption follows social norms \cite{ackert2025bnpl, cardaci2018inequality}, both parameters are evaluated across parameter sweeps.
Control runs set $\beta = 0$ to isolate balance sheet debt-servicing dynamics from peer diffusion channels.
Neither coefficient is fixed by any source: the $\beta = 0$ control arm removes the circularity this would create but not the plausibility problem, since it does not say which region of the $(q_{\text{base}}, \beta)$ surface is the actual market, and the available anchors, the 82.8\% banked ceiling, TransUnion's 57\% holding rate \cite{transunion_cps_sa_2025} and the \ac{CFPB} stacking shares, are indirect.

To address potential sensitivity to linear specifications, we pre-register a Granovetter-style heterogeneous threshold model as a structural robustness check \cite{granovetter1978threshold}.
Under this formulation, household propensity steps to $q_i = q_{\text{base}} + \gamma$ once group adoption $s_{g(i)}$ exceeds an agent-specific threshold $\theta_i \sim \mathcal{N}(\mu_\theta, \sigma_\theta)$, truncated to $[0,1]$.
Truncation retains a subpopulation with zero-threshold values representing early adopters.
Threshold parameters are generated using fixed initialization seeds to maintain exact comparability across control settings.
Parameter sweeps focus primarily on threshold dispersion ($\sigma_\theta$), testing whether cascade dynamics depend on variance in peer susceptibility rather than mean thresholds.

\subsubsection*{Submodel 13: Stacking}
\label{sec:stacking}

Absent centralized reporting, households can maintain concurrent credit lines across multiple platforms without aggregate exposure caps (Submodel~10). We evaluate platform multiplicity by sweeping the available provider count from one to six, treating the single-platform setting as a benchmark to separate multi-platform credit stacking from intra-firm credit expansion. Emerging credit stacking patterns are benchmarked against consumer credit data reported by the \ac{CFPB} (Section~\ref{sec:background}).

\subsubsection*{Submodel 15: Scenario Switches}
\label{sec:levers}

The third research question is answered by running the calibrated model under three scenarios, crossed so that each regulatory switch is observed with and without the behavioural one:
\begin{enumerate}
\item \textbf{Bureau visibility} ($\mathbb{1}^{\text{bur}}$): \ac{BNPL} obligations reach the bureau and enter the Regulation 23A test, closing the asymmetry of Submodel~10; the model's representation of the 2026 reporting requirement \cite{businessday_bnpl_2026}, subject to the limit stated there.
\item \textbf{Mandatory affordability screening} ($\mathbb{1}^{\text{aff}}$): the Regulation 23A test is applied to \ac{BNPL} originations themselves, the South African analogue of the \ac{FCA}'s proportionate affordability check \cite{fca_ps26_1, woolard2021}.
\item \textbf{Peer effects} ($\beta > 0$): adoption is socially transmitted through Equation~\ref{eq:peer}. A behavioural mechanism rather than an instrument, it is a scenario because it asks whether the answers to the other two depend on how adoption spreads; its $\beta = 0$ arm is the control of every comparison.
\end{enumerate}
Two remarks bear on the second scenario.
The \ac{NCR}'s view that a \ac{BNPL} agreement may become an incidental credit agreement once default charges are levied (Section~\ref{sec:background}) implies that the late fees of Submodel~14 could already bring a facility within the \ac{NCA}; the origination screen of Submodel~12 nonetheless stays light, because no provider applies the Regulation 23A test at origination whatever the eventual classification, so the scenario is best read as a clarification of law that may already apply rather than as new regulation.
A scenario's effect is judged on cumulative \ac{BNPL} volume as well as default, so borrowing merely deferred is distinguished from borrowing forgone.
A cooling-off window and a concurrent-facility cap are implemented and swept; their results are reported only in Appendix~\ref{app:supplementary}.

\subsection{Implementation and Experimental Protocol}
\label{sec:implementation}

The model is implemented in Python on the Mesa framework.
Households, the lender, the bureau and the platforms are each a class; reference groups are keys
on the model.
The population loads directly from the parquet files of
Section~\ref{ch:population}, and the Regulation 23A gate reuses the same function that
constructed scheduled servicing, so the gate and the data layer are consistent.
The parameter register of Appendix~\ref{app:specification} is generated from the code's own parameter declarations, so the table and the implementation cannot drift apart.

Model verification is confirmed using automated unit test suites covering balance sheet conservation rules, arrears transition logic across \ac{CCMR} age brackets, statutory Regulation 23A benchmark calculations, pay-in-four repayment execution, and legal late-fee boundaries.
Degenerate parameter sweeps confirm baseline stability: disabling \ac{BNPL} recovers baseline state conditions, setting $\beta = 0$ and $\gamma = 0$ reproduces single-agent behavior, limiting platforms to one eliminates cross-platform stacking, and adjusting payment friction alters arrears distributions without shifting core default rates.

\label{sec:lim-calibration}
Model calibration targets two parameters: the wage-shock onset probability, fitted to the \ac{CCMR} 90+ day share, and the payment-friction rate, fitted to the 1--30 day share; every other parameter is derived from data or swept.
The baseline is therefore calibrated, not validated, and this is the most important limitation in the thesis.
The shock probability cannot be estimated from a single \ac{NIDS} wave, so it is fitted until baseline arrears reproduce Pattern~2, and agreement on the fitted bands is not evidence that the model is right; only the bands left free, and the results obtained after \ac{BNPL} is injected with that parameter held fixed, are out-of-sample tests.
The fitted value is also implausibly high, roughly four times the separation rate the labour force survey records, because the model has exactly one thing that can go wrong where real households also fall behind through illness, unexpected bills, rate rises and family change: one channel does the work of several.
This is disclosed rather than corrected, since correcting it means adding shock channels the data cannot parameterise.
Experiments run as multi-seed Monte Carlo; sensitivity analysis, noise tolerances and the full specification are in Section~\ref{sec:results-robustness}.
